\begin{description}
\item[(A)] The orbit of Earth is elliptical, so the shape of the orbit after
  the solar catastrophe will depend on the moment when the decrees % sic: "decrees" (decrease) in source
  of the mass of the Sun will occur.

  Initial analysis of the problem:
  \begin{description}
  \item[a)] In 3rd July the Earth is at the aphelion. The speed of the Earth
    is smaller than the speed of Earth on a circular orbit with radius
    $r_{0,\mathrm{max}} = a_0(1 + e_0)$.
  \item[b)] In 3rd January the Earth is at perihelion. The speed of the Earth
    is bigger than the speed of Earth on a circular orbit with radius
    $r_{0,\mathrm{max}} = a_0(1 - e_0)$. % sic: source writes r_{0,max} here; should be r_{0,min}
  \end{description}

  Conclusion (A): the period should be calculated only for situation a). The
  expected trajectory in this case is an elliptic one. The possibility that
  Earth hit the Sun is available too.

\item[(B)] Calculations:

  In 3rd July the distance from Sun is maximum: fig.~2.1,
  \[ r_{0,\mathrm{max}} = a_0(1 + e_0). \]
% FIGURE NEEDED: Figure 2.1 -- initial (small, dashed) Earth orbit around S_0; M_0 with v_0,per and v_0,aph marked, and the large final elliptical Earth orbit around S; M with r_min and v_per, showing r_0,max = r_0,aph = r_per of the new orbit (2014 TH solutions, p. 5)

  \textbf{Before the catastrophe:}
  \begin{itemize}
  \item $v_{0,\mathrm{aph}}$ -- the speed of Earth on aphelion,
  \item $a_0$ -- big Earth's elliptical orbit semi axis,
  \item $v_0$ -- the speed of Earth if its orbit is circular with radius
    $r_0 = a_0$.
  \end{itemize}

  According to Keppler's % sic: "Keppler's" in source
  second law and the law of energy conservation (see figure 2.1) the
  following relations can be written:
  \begin{align*}
    & v_{0,\mathrm{per}}\, r_{0,\mathrm{per}}
      = v_{0,\mathrm{aph}}\, r_{0,\mathrm{aph}}; \\
    & \frac{v_{0,\mathrm{per}}^2}{2} - K\frac{M_0}{r_{0,\mathrm{per}}}
      = \frac{v_{0,\mathrm{aph}}^2}{2} - K\frac{M_0}{r_{0,\mathrm{aph}}} \\
    & r_{0,\mathrm{min}} = r_{0,\mathrm{per}} = a_0(1 - e_0) \\
    & r_{0,\mathrm{max}} = r_{0,\mathrm{aph}} = a_0(1 + e_0) \\
    & KM_0 = v_0^2\, r_0 = v_0^2\, a_0 \\
    & v_0 = \sqrt{K\frac{M_0}{r_0}} = \sqrt{K\frac{M_0}{a_0}} \\
    & v_{0,\mathrm{per}} = v_0 \sqrt{\frac{1 + e_0}{1 - e_0}} > v_0 \\
    & v_{0,\mathrm{per}} > v_0 \tag*{(1)} \\
    & v_{0,\mathrm{aph}} = v_0 \sqrt{\frac{1 - e_0}{1 + e_0}} \\
    & v_{0,\mathrm{aph}} < v_0 \tag*{(2)}
  \end{align*}

  Conclusion -- According to the relations (1) and (2) the new orbit of the
  Earth could be an elliptic one.

  For the new elliptical Earth orbit:
  \begin{align*}
    & r_{\mathrm{per}} = r_{0,\mathrm{aph}}; \\
    & r_{\mathrm{min}} = r_{\mathrm{per}} = a(1 - e); \\
    & a_0(1 + e_0) = a(1 - e); \quad a = a_0\,\frac{1 + e_0}{1 - e}; \\
    & v_{\mathrm{per}} = v_{0,\mathrm{aph}}; \\
    & v_{\mathrm{per}} = v\sqrt{\frac{1 + e}{1 - e}},
  \end{align*}
  where $v$ este % sic: Romanian "este" (is) left in source
  the Earth's speed on a circular orbit with the radius $r = a$, when the
  mass of the Sun becomes $M = M_0/2$;
  \begin{align*}
    & v\sqrt{\frac{1 + e}{1 - e}} = v_0\sqrt{\frac{1 - e_0}{1 + e_0}}; \\
    & v = \sqrt{K\frac{M}{r}} = \sqrt{K\frac{M_0}{2a}}
      = \sqrt{K\frac{M_0}{a_0}}\sqrt{\frac{a_0}{2a}}
      = v_0\sqrt{\frac{a_0}{2a}}; \\
    & e = 1 - 2e_0; \quad a = a_0\,\frac{1 + e_0}{2e_0}.
  \end{align*}

  Conclusion:
  \begin{align*}
    T_0 &= \frac{2\pi r_0}{v_0} = \frac{2\pi a_0}{v_0} \\
    T &= \frac{2\pi r}{v} = \frac{2\pi a}{v}; \\
    \frac{T}{T_0} &= \frac{a}{a_0}\,\frac{v_0}{v}
      = \frac{1 + e_0}{2e_0}\sqrt{\frac{2a}{a_0}}
      = \frac{1 + e_0}{2e_0}\,\sqrt{2}\,\sqrt{\frac{1 + e_0}{2e_0}}; \\
    T &= T_0\sqrt{2}\left(\frac{1 + e_0}{2e_0}\right)^{3/2}
      \approx 230\ \text{years}
  \end{align*}

  b) In 3rd of January the Earth is at perihelion. In that moment the Erath % sic: "Erath" in source
  speed is larger than the speed necessary for an Earth's circular orbit.
  Thus the trajectory of the Earth after the catastrophe will be an open
  trajectory, i.e.\ an hyperbolic or parabolic orbit.

  Conclusion: it is not necessary to calculate the period of revolution or it
  could be issued as infinite.
\end{description}
